\section{Búsqueda, Critique y Preference Learning: el mecanismo de entrenamiento de AgentQ}

El artículo de AgentQ combina tres mecanismos: Monte Carlo Tree Search (MCTS) amplía la exploración, la LLM critique ofrece retroalimentación de proceso sobre los nodos intermedios y un objetivo al estilo de Direct Preference Optimization (DPO) aprende de las trayectorias mejores y peores. Los tres resuelven, respectivamente, los problemas de «no hay exploración», «no se sabe dónde está el error» y «la experiencia exitosa no se consolida».

\lecturefigure{slide-34-agentq-methods-paper.jpg}{Los tres componentes de AgentQ: MCTS, self-critique/process supervision y DPO}{00:20:47}

\subsection{MCTS: conservar ramas alternativas para tareas largas}

La búsqueda en árbol conserva varias acciones candidatas en un mismo estado, en lugar de comprometerse de forma voraz con la primera ruta. La fórmula de selección UCT típica es

Este paso retoma la autocorrección vista en la demo: si la policy emite una sola acción cada vez, después de un fallo solo puede volver a muestrear desde el mismo contexto, y es fácil que repita la ruta original; la estructura de árbol, en cambio, guarda de forma explícita «qué ya se probó, qué ramas tienen mayor valor y dónde todavía vale la pena explorar». Al leer una figura de MCTS, primero se identifican la root, las ramas candidatas y los nodos hoja, y después se distingue entre la value estimate y el resultado real de la ejecución. Para un agent web, el árbol de búsqueda debe desplegarse de preferencia en un simulador o en una copia que admita reversión; en el sitio real solo se ejecuta un número reducido de acciones ya filtradas.
\[
a^*=\arg\max_a\left[Q(s,a)+c\sqrt{\frac{\ln N(s)}{N(s,a)+1}}\right].
\]
$Q(s,a)$ es la estimación de valor actual, $N(s)$ es el número de visitas al estado padre, $N(s,a)$ es el número de visitas a la acción y $c$ controla la exploration. El primer término favorece las rutas que ya se sabe que son buenas; el segundo incentiva probar las acciones menos visitadas.

\lecturefigure{slide-35-mcts-search.jpg}{MCTS equilibra exploración y explotación en el árbol de decisiones del Agent}{00:21:40}

\begin{warningbox}{MCTS en entornos web no es un retroceso gratuito}
Muchas acciones son irreversibles: enviar un pedido, mandar un mensaje o realizar un pago no se pueden deshacer como en un tablero de ajedrez. Una implementación segura debe buscar en un entorno simulado o desplegar solo estados que admitan reversión; la exploración en el entorno real debe tener límites de permisos, de presupuesto y de sandbox.
\end{warningbox}

\subsection{Process feedback: recompensar las decisiones intermedias, no solo el desenlace}

Usar únicamente el éxito final $r_{\text{outcome}}$ produce una señal dispersa y, además, puede etiquetar como ejemplo positivo una trayectoria que «tuvo éxito por casualidad, pero con un proceso peligroso». La recompensa de proceso puede escribirse como
\[
R(\tau)=\alpha r_{\text{outcome}}+
\sum_{t=1}^{T}\beta_t r_{\text{process},t}-\lambda C(\tau),
\]
$\tau$ es la trayectoria completa, $r_{\text{process},t}$ evalúa si el paso $t$ es razonable y $C(\tau)$ registra el costo de las herramientas, las infracciones o el riesgo irreversible. $\alpha,\beta_t,\lambda$ determinan el equilibrio entre resultado, proceso y costo.

\lecturefigure{slide-36-self-critique-process.jpg}{La LLM critique ofrece retroalimentación de proceso sobre las acciones candidatas}{00:22:58}

\begin{knowledgebox}{Los tres niveles de la critique}
La \textbf{syntax critique} revisa el formato de la acción; la \textbf{state critique} revisa si la acción es coherente con la página actual y sus restricciones; la \textbf{goal critique} revisa si la acción sigue sirviendo al objetivo original del usuario. Una critique que solo evalúa la fluidez del lenguaje no puede detectar errores de fecha, identidad, permisos ni presupuesto.
\end{knowledgebox}

\subsection{Construir preferencias a partir de ramas exitosas y fallidas}

AgentQ ordena las ramas que produce la búsqueda en preferred y rejected, y después actualiza la policy con un objetivo de preferencias fuera de línea. La forma habitual de DPO es
\[
\mathcal{L}_{\text{DPO}}(\theta)
=-\mathbb{E}\log\sigma\!\left(
\beta\left[
\log\frac{\pi_\theta(y^+\mid x)}{\pi_{\text{ref}}(y^+\mid x)}-
\log\frac{\pi_\theta(y^-\mid x)}{\pi_{\text{ref}}(y^-\mid x)}
\right]\right).
\]
$x$ es el estado junto con el contexto de la tarea, $y^+$ es la acción o trayectoria mejor, $y^-$ es la rama peor, $\pi_{\text{ref}}$ es la política de referencia y $\beta$ controla la intensidad de la preferencia. En el caso de los agents, las preferencias no pueden provenir solo de la captura de pantalla final de la página; deben combinarse con la critique de proceso, los errores de las herramientas y las restricciones de seguridad.

\lecturefigure{slide-37-rlhf-preference-learning.jpg}{Construir señales de aprendizaje de preferencias a partir de trayectorias de búsqueda y critique}{00:23:23}

\teachervoice{00:20:47--00:23:24, el ponente explica AgentQ mediante MCTS, LLM critique y preference learning. El docente enfatiza: tanto las trayectorias exitosas como las fallidas deben entrar en el aprendizaje; sobre todo, hay que lograr que el modelo aprenda a hacer backtrack y recover en entornos complejos.}

\subsection{La trayectoria de OpenTable en siete pasos: el éxito final no oculta los errores intermedios}

Este conjunto de figuras debe leerse como una secuencia de transiciones de estado, no como siete capturas parecidas. La trayectoria primero navega hacia el restaurante objetivo, entra a la página principal, localiza el restaurante y luego selecciona una fecha equivocada; después abre el selector de fechas, corrige la fecha y completa la reservación. El error ocurre en el cuarto paso, la recuperación en el quinto y el sexto, y la confirmación final en el séptimo.

El valor didáctico de esta trayectoria es que incluye a la vez la \textbf{conservación del objetivo} y la \textbf{corrección local}. El sistema no se puso a buscar otro restaurante porque la fecha estuviera mal, ni redujo la solicitud del usuario a «basta con reservar en cualquier restaurante»; conservó el restaurante, el número de personas y el rango de horario, y solo modificó el campo erróneo. Al leer cada figura conviene seguir tres líneas: el estado actual de la página, la siguiente acción que el agent declara y las restricciones del objetivo original que todavía no se cumplen. Solo así se puede juzgar si la recuperación es una corrección real o si se llegó por casualidad a un resultado que podía completarse.

\lecturefigure{slide-38-opentable-step1-navigate.jpg}{Paso 1: navegar al restaurante y al sitio objetivo}{00:24:10}

\lecturefigure{slide-39-opentable-step2-homepage.jpg}{Paso 2: entrar a la página principal de OpenTable e iniciar la búsqueda}{00:24:15}

\lecturefigure{slide-40-opentable-step3-restaurant.jpg}{Paso 3: localizar la página de La Pizza \& La Pasta}{00:24:42}

\lecturefigure{slide-41-opentable-step4-wrong-date.jpg}{Paso 4: se seleccionó una fecha equivocada; la trayectoria presenta un fallo diagnosticable}{00:24:49}

\lecturefigure{slide-42-opentable-step5-date-selector.jpg}{Paso 5: volver a abrir el selector de fechas}{00:24:52}

\lecturefigure{slide-43-opentable-step6-correct-date.jpg}{Paso 6: corregir la fecha y entrar al flujo de confirmación}{00:24:54}

\lecturefigure{slide-44-opentable-step7-confirmed.jpg}{Paso 7: la reservación se completa; la página de confirmación aporta la evidencia final}{00:24:56}

\begin{importantbox}{Trajectory-level acceptance}
Registrar el éxito final como $s_T=1$ todavía no basta. Se recomienda reportar al mismo tiempo: el número de acciones erróneas, el número de recuperaciones, el número de acciones irreversibles, el número de confirmaciones del usuario, el costo total de las herramientas y la latencia hasta completar la tarea. Solo así se distingue entre «completarse sin contratiempos a la primera» y «dar rodeos, hacer clics erróneos, pero acabar teniendo éxito por casualidad».
\end{importantbox}

\begin{lstlisting}[caption={Pseudocódigo de verificación y recuperación de trayectorias}]
for step in range(max_steps):
    observation = environment.observe()
    action = policy.propose(goal, observation, memory)
    if not permission_check(action):
        action = request_user_confirmation(action)
    result = environment.execute(action, idempotency_key=step)
    if postcondition_holds(goal, result):
        return SUCCESS
    if repeated_state(result) or budget_exhausted():
        return HUMAN_FALLBACK
    memory.record(observation, action, result)
return TIMEOUT
\end{lstlisting}

\subsection{Página de resultados: las cifras deben ir ligadas a la configuración experimental}

La gráfica de tasa de éxito en OpenTable que se mostró en clase compara varias configuraciones y destaca la mejora que aportan componentes como AgentQ y online search. El artículo reporta que Llama-3 70B sin ejemplos pasa de 18.6\% a 81.7\%, y que al agregar online search alcanza 95.4\%. Estas cifras son resultados bajo una tarea, una recolección de datos, un modelo y un método de evaluación específicos; respaldan que «la búsqueda y el aprendizaje son eficaces», pero no que «todas las tareas web ya son confiables al 95\%».

Al leer la gráfica de barras, primero se compara el cambio al agregar componentes sobre el mismo base model y después las diferencias entre modelos distintos; no hay que atribuir a un solo algoritmo la diferencia entre todas las barras. También hay que revisar si el criterio de éxito considera solo la página final, si se permite búsqueda humana o externa, cuántos datos de interacción usó cada configuración y si los fallos provienen de la misma distribución de tareas. Si las configuraciones experimentales difieren, la altura de las barras solo puede tomarse como el resultado de la combinación del sistema; de ella no se puede deducir directamente la contribución causal independiente de un solo módulo de MCTS, DPO o search.

\lecturefigure{slide-45-agentq-real-world-benchmark.jpg}{Comparación de la tasa de éxito de AgentQ en escenarios reales de reservación}{00:26:41}

\begin{warningbox}{No hay que confundir la mejora relativa con la confiabilidad absoluta}
Pasar de 18.6\% a 81.7\% es un enorme avance de investigación, pero 81.7\% significa que aproximadamente una de cada cinco tareas sigue fallando. Incluso con 95.4\%, si se ejecutan cien acciones de alto riesgo al día, el número esperado de fallos sigue siendo inaceptable. El umbral para llevarlo a producción debe considerar la frecuencia, la gravedad y la posibilidad de recuperación.
\end{warningbox}

\subsection{Resumen de la sección}

La cadena de mecanismos de AgentQ es la siguiente: la búsqueda en árbol genera ramas diversas, la process critique evalúa los pasos intermedios, el objetivo de preferencias reincorpora a la política las ramas buenas y malas, y la evaluación a nivel de trayectoria revisa a la vez el resultado y el proceso. El caso de OpenTable muestra que los errores no son lo preocupante; los defectos sistémicos son no poder detectarlos, no poder recuperarse de ellos y no poder aprender de ellos.
